% Typeset using lualatex

\documentclass[a4paper,10pt]{article}
\usepackage[utf8]{luainputenc}
\usepackage[pdfborder={0 0 0}]{hyperref}

\usepackage[czech]{babel}

\usepackage{fouriernc}
\usepackage{tgpagella}
\usepackage{xcolor}
\usepackage[activate=true]{microtype}
\usepackage{amssymb}
\usepackage{fontspec}
\usepackage{titling}
\usepackage{titlesec}

\usepackage{calc}
\usepackage{multicol}
\usepackage{array}
\usepackage{enumitem}
\usepackage{tikz}
\usepackage{siunitx}
\usepackage{colortbl}
\usepackage{everypage}
\usepackage[a4paper,margin=2cm,top=1cm]{geometry}

\setmainfont[Numbers=OldStyle]{TeX Gyre Pagella}

\definecolor{plpink}{HTML}{FF3232}
\definecolor{plhome}{HTML}{CCEEFF}
\definecolor{silver}{HTML}{888888}

\sisetup{
    output-decimal-marker={,}
}
\pagestyle{empty}

\parskip=1ex
\parindent=0pt

\setlist[enumerate,1]{start=0}

\AddEverypageHook{
    \begin{tikzpicture}[remember picture,overlay]
        \node[opacity=0.1,below left,xshift=1cm,yshift=1cm] at (current page.north east) {\includegraphics[width=12cm]{../../images/pylady-pink}};
        \node[below left,align=right,yshift=-1.5em] at (current page.north east) {\color{white} strana \thepage};
        \node[below left,align=right] at (current page.north east) {\color{white} sada \plsetno};
    \end{tikzpicture}
}

\newcommand\answerspace{\\\rule[0cm]{0pt}{1cm}}

\newcommand\True{\texttt{True}}
\newcommand\False{\texttt{False}}

\usetikzlibrary{turtle}
\newcommand\turtle[1]{\tikz[>=stealth,x=0.25mm,y=0.25mm]\draw[turtle={home,rt=90,#1}];}
\newcommand\gulp[1]{}

\newcommand\plsetno{5}

\newcommand\startsection[1]{
     \vspace{0.2ex}
    \hrule
    {\fontspec{Oxygen} \tiny
     \vspace{-1ex}
     \emph{#1}
     \vspace{-1.5em}
    }
}

\newfontfamily\headingfont[]{Bree Serif}
\titleformat*{\section}{\LARGE\headingfont}

\begin{document}

\section*{Domácí úkoly \plsetno}

Nezapomeň: Když něco funguje, přidej to do Gitu!
\\\emph{\small Pokud materiály procházíš v jiném pořadí než PyLadies Brno 2015, a Git ještě neznáš, můžeš všechny poznámky o Gitu ignorovat.}

\vspace{1cm}

\startsection{Trocha experimentování. Zkus se zamyslet, jestli jsi „dobře” pochopila otázku.}

\begin{enumerate}[resume]

\item Co se stane, když tělo nějaké funkce necháme prázdné?

\item Co se stane, když necháme prázdné tělo cyklu?

\end{enumerate}

\startsection{Procvičení funkcí. Jestli jsi Pythoní funkce nepsala už před kurzem, tak první z těchto úkolů určitě udělej. Druhý jen pokud máš ráda geometrii :)}

\begin{enumerate}[resume]

\item Napiš funkci, která vykreslí domeček dané velikosti.
    \\\emph{\small (t.j. velikost se zadá argumentem)}

\gulp{
from math import sqrt
from texturtle import ctx, forward, left, right

def domecek(velikost):
    for i in range(4):
        forward(velikost)
        left(90)
    left(45)
    for n in 1, 2, 2, 1:
        forward(velikost*sqrt(2)/n)
        left(90)
    right(45)

with ctx():
    for n in 10, 20, 50:
        domecek(n)
        forward(n + 10)
}
\turtle{fd=10,lt=90,fd=10,lt=90,fd=10,lt=90,fd=10,lt=90,lt=45,fd=14.142135623730951,lt=90,fd=7.0710678118654755,lt=90,fd=7.0710678118654755,lt=90,fd=14.142135623730951,lt=90,rt=45,fd=20,fd=20,lt=90,fd=20,lt=90,fd=20,lt=90,fd=20,lt=90,lt=45,fd=28.284271247461902,lt=90,fd=14.142135623730951,lt=90,fd=14.142135623730951,lt=90,fd=28.284271247461902,lt=90,rt=45,fd=30,fd=50,lt=90,fd=50,lt=90,fd=50,lt=90,fd=50,lt=90,lt=45,fd=70.71067811865476,lt=90,fd=35.35533905932738,lt=90,fd=35.35533905932738,lt=90,fd=70.71067811865476,lt=90,rt=45,fd=60,}
    \\\emph{\small Funguje? Do Gitu s tím!}

\item Máš-li ráda geometrii, můžeš zkusit dávat domečkové funkci dva argumenty:
    šířku a výšku.
    \\\emph{\small
        Je potřeba si vzpomenout na Pythagorovu větu a funkci tangens.
        Pozor, funkce \texttt{tan} vrací výsletek v radiánech;
        je potřeba ho převést na stupně.
    }

\gulp{
from math import sqrt, degrees, atan
from texturtle import ctx, forward, left, right

def domecek(sirka, vyska):
    for i in range(2):
        forward(sirka)
        left(90)
        forward(vyska)
        left(90)

    uhel = degrees(atan(vyska / sirka))
    uhlopricka = sqrt(sirka**2 + vyska**2)
    left(uhel)
    forward(uhlopricka)
    left(180 - 2 * uhel)
    forward(uhlopricka / 2)
    left(2 * uhel)
    forward(uhlopricka / 2)
    left(180 - 2 * uhel)
    forward(uhlopricka)
    left(uhel)

with ctx():
    for a, b in (20, 20), (10, 50), (60, 20):
        domecek(a, b)
        forward(a + 10)
}
\turtle{fd=20,lt=90,fd=20,lt=90,fd=20,lt=90,fd=20,lt=90,lt=45.0,fd=28.284271247461902,lt=90.0,fd=14.142135623730951,lt=90.0,fd=14.142135623730951,lt=90.0,fd=28.284271247461902,lt=45.0,fd=30,fd=10,lt=90,fd=50,lt=90,fd=10,lt=90,fd=50,lt=90,lt=78.69006752597979,fd=50.99019513592785,lt=22.61986494804043,fd=25.495097567963924,lt=157.38013505195957,fd=25.495097567963924,lt=22.61986494804043,fd=50.99019513592785,lt=78.69006752597979,fd=20,fd=60,lt=90,fd=20,lt=90,fd=60,lt=90,fd=20,lt=90,lt=18.43494882292201,fd=63.245553203367585,lt=143.13010235415598,fd=31.622776601683793,lt=36.86989764584402,fd=31.622776601683793,lt=143.13010235415598,fd=63.245553203367585,lt=18.43494882292201,fd=70,}

\end{enumerate}

\startsection{Procvičení programování. Často je jednoduché něco napsat, ale dotažení do konce může být časově náročné. Nemáš-li čas, zkus se aspoň zamyslet jak bys úkol vyřešila.}

\begin{enumerate}[resume]

\item Změň program Kámen, Nůžky, Papír tak, aby opakoval hru dokud
    uživatel nezadá \texttt{"konec"}.

\item Změň funkci \texttt{ano\_nebo\_ne} tak,
    aby se místo \texttt{"ano"} se dalo použít i \texttt{"a"},
    místo \texttt{"ne"} i \texttt{"n"},
    a aby se nebral ohled na velikost písmen a mezery před/za odpovědí.
    \\\emph{\small
        Textům jako \texttt{"možná"} nebo \texttt{"no tak určitě"}
        by počítač dál neměl rozumět.
    }

\end{enumerate}

\startsection{Sada zajímavých (snad) prográmků, které bys teď měla být schopná napsat. Nemáš-li čas, zatím je přeskoč.}

\begin{enumerate}[resume]

\item Napiš program, který se zeptá na příjmení uživatelky/uživatele,
    a zkusí podle něj uhodnout její/jeho pohlaví.
    \\\emph{\small Připomínám: Až to bude fungovat, dej to do Gitu!}

\item Najdi na internetu text své oblíbené písně, zkopíruj si ho do řetězce,
    a zjisti, kolikrát je v něm použito písmeno K.
    \\\emph{\small Připomínám: Až to bude fungovat, dej to do Gitu!}

\item Napiš program, který simuluje tuto hru:
    \\První hráč hází kostkou
    \emph{(t.j. vybírají se náhodná čísla od 1 do 6)},
    dokud nepadne šestka.
    Potom hází další hráč, dokud nepadne šestka i jemu.
    Potom hází hráč třetí, a nakonec čtvrtý.
    Vyhrává ten, kdo na hození šestky potřeboval nejvíc hodů.
    (V případě shody vyhraje ten, kdo házel dřív.)
    \\Program by měl vypisovat všechny hody, a nakonec napsat, kdo vyhrál.
    \\\emph{\small Připomínám: Až to bude fungovat, dej to do Gitu!}

\end{enumerate}

\vfill

Pokračování na další straně. \hfill \emph{\%}

\newpage

\startsection{Nakonec trošku delší projekt. Budeme na něm stavět dál; nedokončíš-li ho teď, budeš ho muset dodělat před příští sadou úkolů.}

\begin{center}
\fbox{
\begin{minipage}{8cm}
1-D piškvorky se hrají na řádku s dvaceti políčky.
    \\Hráči střídavě přidávají kolečka (\texttt{o}) a křížky (\texttt{x}),
        třeba:
    \\1. kolo: \texttt{-------\textcolor{plpink}{x}------------}
    \\2. kolo: \texttt{-------x--\textcolor{plpink}{o}---------}
    \\3. kolo: \texttt{-------x\textcolor{plpink}{x}-o---------}
    \\4. kolo: \texttt{-------xx\textcolor{plpink}{o}o---------}
    \\5. kolo: \texttt{------\textcolor{plpink}{x}xxoo---------}
    \\Hráč, která dá tři své symboly vedle sebe, vyhrál.
\end{minipage}
}
\end{center}

\begin{enumerate}[resume]

\item Napiš funkci \texttt{vyhodnot}, která dostane řetězec
    s herním polem 1-D piškvorek,
    a vrátí jednoznakový řetězec podle stavu hry:
    \begin{enumerate}
    \item[\texttt{"x"}] – Vyhrál hráč s křížky (pole obsahuje \texttt{xxx})
    \item[\texttt{"o"}] – Vyhrál hráč s kolečky (pole obsahuje \texttt{ooo})
    \item[\texttt{"!"}] – Remíza (pole neobsahuje \texttt{-}, a nikdo nevyhrál)
    \item[\texttt{"-"}] – Ani jedna ze situací výše
    \end{enumerate}
    \emph{\small Připomínám: Až to bude fungovat, dej to do Gitu!}

\item Napiš funkci \texttt{tah}, která dostane řetězec s herním polem,
    číslo políčka (0-19), a symbol (\texttt{x} nebo \texttt{o}),
    a vrátí herní pole \emph{(t.j. řetězec)} s daným symbolem umístěným na danou pozici.

    Hlavička funkce by tedy měla vypadat nějak takhle:
\\\verb+    def tah(pole, cislo_policka, symbol):+
\\\verb+        "Vrátí herní pole s daným symbolem umístěným na danou pozici"+
\\\verb+        ...+
    \\\emph{Můžeš využít nějakou funkci, kterou jsme napsaly už na sraze?}

\item Napiš funkci \texttt{tah\_hrace}, která dostane řetězec s herním polem,
    zeptá se hráče, na kterou pozici chce hrát, a vrátí herní pole
    se zaznamenaným tahem hráče.
    \\Funkce by měla odmítnout záporná nebo příliš velká čísla,
    a tahy na obsazená políčka.
    Pokud uživatel zadá špatný vstup, funkce mu vynadá a zeptá se znova.
    \\\emph{\small Funguje? Do Gitu s tím!}

\item Napiš funkci \texttt{tah\_pocitace}, která dostane řetězec s herním polem,
    vybere pozici, na kterou hrát, a vrátí herní pole
    se zaznamenaným tahem počítače.
    \\Použij jednoduchou náhodnou „strategii”:
    \begin{enumerate}
    \item[1.] Vyber číslo od 0 do 19
    \item[2.] Pokud je dané políčko volné, hrej na něj
    \item[3.] Pokud ne, opakuj od bodu 1
    \end{enumerate}

    Hlavička funkce by tedy měla vypadat nějak takhle:
\\\verb+    def tah_pocitace(pole):+
\\\verb+        "Vrátí herní pole se zaznamenaným tahem počítače"+
\\\verb+        ...+

\item Napiš funkci \texttt{piskvorky1d}, která vytvoří řetězec s herním polem,
    a střídavě volá funkce \texttt{tah\_hrace} a \texttt{tah\_pocitace},
    dokud někdo nevyhraje nebo nedojde k remíze.
    \\Nezapomeň kontrolovat stav hry po každém tahu.
    \\\emph{\small Funguje? Do Gitu s tím!}

\end{enumerate}

\startsection{Poslední úkol je nepovinný, ale, jak to u podobných úkolů bývá, můžeš na něj dostat zpětnou vazbu. Doporučuju toho využít!}

\begin{enumerate}[resume]

\item Zvládneš pro počítač naprogramovat lepší strategii?
    Třeba aby se snažil hrát vedle svých existujících symbolů,
    nebo aby bránil protihráčovi?
    \\Pokud ano, pošli svoji funkci \texttt{tah\_pocitace}
    na soukromý e-mail organizátora (ne do diskusní skupiny).
    Pošli ho jako přílohu, nekopíruj ho do textu e-mailu.
    \\Stačí jen docela malé vylepšení!
    \\\emph{\small A do Gitu to samozřejmě dej taky...}

\end{enumerate}

\end{document}

